\reviewer

\begin{generalcomment}
The revised manuscript is substantially clearer and is recommended for acceptance after minor revision. The remaining comments concern the hash-length convention, probability wording, provisional deployment guidance, reproducibility materials, and the interpretation of qualified security-table entries.
\end{generalcomment}

\begin{revcomment}
Section~6.2.1 uses 160-bit hashes and fields for analytical accounting, whereas the host benchmark uses SHA-256. Please clarify whether SHA-256 outputs are truncated or whether a different field-width configuration is used. Explain the consequence for Table~4 and the challenge-expansion cost in Section~3.1.1.
\end{revcomment}

\begin{revresponse}
We agree that the previous wording left the hash computation width and the serialized field width implicit. The revised manuscript now defines $h(m)=\operatorname{Trunc}_{160}(\operatorname{SHA256}(m))$ for protocol fields. Thus, the 160-bit profile used by Table~4 is an explicit protocol configuration, not an unqualified claim that SHA-256 produces a 160-bit digest. The host benchmark measures the underlying full SHA-256 operation; truncation is a serialization step after the compression operation. Each challenge-expansion call therefore contributes a 20-byte output, and the expected $2.000008T_H$ expansion cost is already included in the LD and EN hash counts. A deployment using full 256-bit fields would require separate communication and storage accounting. The benchmark script now exposes the 160-bit field conversion while retaining full-SHA-256 timing for the primitive measurement.
\end{revresponse}
\begin{changes}
\textbf{Manuscript change.} The hash-length convention was added to the Preliminaries and challenge-expansion discussion. The Table~4 caption and cost model now identify the 160-bit field profile explicitly, and the benchmark script documents the distinction between full SHA-256 timing and 160-bit protocol-field serialization.
\end{changes}

\begin{revcomment}
Table~3 footnote b describes $q_c/26^s$ as a success probability per admitted trial. Please distinguish the single-guess probability, the union bound over admitted trials, and the ROR advantage term.
\end{revcomment}

\begin{revresponse}
We corrected this wording. A single admitted CEAF guess succeeds with probability $26^{-s}$. Across $q_c$ admitted online trials, the success probability is bounded by $\min\{1,q_c/26^s\}$. The corresponding ROR advantage contribution is $2\min\{1,q_c/26^s\}$; the complete reduction additionally includes offline tests through $q_e+q_c$. The revised Table~3 note now states these three quantities separately.
\end{revresponse}
\begin{changes}
\textbf{Manuscript change.} Table~3 footnote b was rewritten, and the probability terminology was checked against the ROR reduction and the sensitivity discussion.
\end{changes}

\begin{revcomment}
The description of $s=3$ as a recommended general setting is stronger than the analytical evidence because usability testing remains future work.
\end{revcomment}

\begin{revresponse}
We agree and have softened the recommendation. Table~5 now labels its final column as provisional guidance. The $s=3$ row is described as a candidate reference configuration whose suitability depends on the admitted-attempt budget, observation environment, and operational workflow. The manuscript does not present this table as empirical usability evidence and retains the planned usability validation as future work.
\end{revresponse}
\begin{changes}
\textbf{Manuscript change.} Table~5, the CEAF parameter discussion, and the usability response were revised to replace ``recommended general setting'' with provisional, workflow-dependent guidance.
\end{changes}

\begin{revcomment}
The benchmark script and HLPSL artifacts should be identified by supplementary-file location or a stable repository link. A brief description of the fuzzy-reproduction routine would also improve reproducibility.
\end{revcomment}

\begin{revresponse}
We have added a reproducibility manifest to the Data Availability Statement and supplied the artifacts in the public \href{https://github.com/secagg4AsyncFL/A-Lightweight-Four-Credential-Authentication-and-Key-Agreement-Protocol-for-AI-Enabled-Edge-IIoT}{GitHub repository}. The benchmark script and recorded JSON results are under \texttt{benchmarks/}; the HLPSL model, AVISPA outputs, and WSL instructions are under \texttt{avispa/}. The reference routine uses a deterministic 256-byte biometric sample, a 16-byte salt, a SHA-256 helper digest, and two HMAC-SHA-256 calls. It is explicitly provided as a timing surrogate rather than as a mandated error-correcting fuzzy extractor.
\end{revresponse}
\begin{changes}
\textbf{Manuscript change.} The Data Availability Statement now lists Supplementary Files S1--S6 and describes the reference fuzzy-reproduction routine. The same files and execution instructions have been synchronized to the public GitHub repository.
\end{changes}

\begin{revcomment}
Qualified entries such as sensor-capture resistance and tracking resistance may appear unconditional when read without the table notes. Please use ``Partial'' or an equivalent explicit label.
\end{revcomment}

\begin{revresponse}
We agree. Table~3 now defines $P$ as a partial or qualified property, alongside $Y$ and $N$. The qualified entries for the proposed scheme are marked with $P$ and the corresponding superscript. In particular, sensor-capture resistance is marked partial because the captured node itself remains compromised while lateral compromise of uncaptured nodes is contained; tracking resistance is marked partial because the claim is passive unlinkability and an active drop may link adjacent runs. The same notation is applied to the other qualified entries of the proposed scheme for consistency.
\end{revresponse}
\begin{changes}
\textbf{Manuscript change.} Table~3 now uses $P^a$--$P^h$ for qualified properties of the proposed scheme and defines $P$ directly below the table. The notes explicitly state the scope of sensor-capture containment, passive unlinkability, and bounded recovery.
\end{changes}
